\chapter{Implementação em Fortran}
\label{cap:implementacao}

\section{Princípio norteador: extração literal, não reformulação}

A convenção adotada em todo o desenvolvimento -- já estabelecida em
partes anteriores do projeto \selfgrib{} e mantida aqui -- foi a de
\textbf{extração literal} das fórmulas e algoritmos do código-fonte de
referência do MPAS-Model, em vez de reimplementá-los ``do zero'' a partir
da descrição em artigos ou manuais. A justificativa é dupla: primeiro,
garante consistência numérica exata com o que o modelo real produz,
eliminando uma classe inteira de divergência que surgiria de decisões de
implementação sutilmente diferentes (ordem de operações, escolha entre
formas matematicamente equivalentes de uma mesma expressão, tratamento de
casos de borda); segundo, torna a origem de cada linha de código
rastreável e auditável -- cada módulo novo carrega, em comentário de
cabeçalho, a referência exata do arquivo/rotina/faixa de linhas do
código-fonte original de onde foi extraído, e a versão específica
consultada (\texttt{mpas-bundle-3.0.2}, confirmada como o fonte exato que
gerou os binários de produção usados neste projeto, distinta e por vezes
divergente da \emph{branch} principal do repositório público -- por
exemplo, a fórmula do fator de suavização $s(i)$ na
\cref{eq:suavizacao_hx} usa $z_t$ como referência de altura na versão
3.0.2, mas $z_h$ na \emph{branch} \texttt{master} mais recente).

\section{Estrutura de módulos}

A \cref{tab:modulos} resume os módulos Fortran novos, organizados pela
fase a que pertencem, com suas dependências e origem no código-fonte de
referência.

\begin{table}[H]
\centering
\small
\caption{Módulos Fortran novos desenvolvidos neste trabalho.}
\label{tab:modulos}
\begin{tabular}{@{}p{3.1cm}p{1.3cm}p{6.7cm}@{}}
\toprule
Módulo/programa & Fase & Origem no código-fonte de referência \\
\midrule
\texttt{hinterp\_native.F90} & 1 & Cópia estrutural de \texttt{convert\_mpas.F90}, reutilizando sem modificação \texttt{mpas\_mesh}, \texttt{target\_mesh}, \texttt{remapper} \\
\texttt{vertical\_grid.F90} & 2 & \texttt{init\_atm\_case\_gfs}, bloco \texttt{config\_tc\_vertical\_grid} \\
\texttt{vinterp\_native.F90} & 3 & Função local \texttt{vertical\_interp} de \texttt{init\_atm\_case\_gfs} \\
\texttt{hydrostatic.F90} & 4 & \texttt{init\_atm\_case\_gfs}, bloco \texttt{config\_met\_interp} \\
\texttt{surface\_fields.F90} & 6 & \texttt{mpas\_atmphys\_initialize\_real.F}, \texttt{mpas\_atmphys\_date\_time.F} \\
\texttt{gen\_init\_native.F90} & 2+3+4+6+7 & Programa-driver, novo (orquestração + escrita NetCDF) \\
\texttt{gen\_lbc\_native.F90} & 3+4+7 & Programa-driver, novo (reusa grade vertical do \texttt{init.nc}) \\
\texttt{namelist\_config.F90} & todas & Novo (leitura de \textit{namelist} real, sem equivalente direto) \\
\bottomrule
\end{tabular}
\end{table}

\section{Convenção de precisão numérica: uma exceção deliberada}

Todo o restante do projeto \selfgrib{} usa precisão dupla
(\texttt{RKIND}, \texttt{real*8}) como padrão. A única exceção
deliberada é o bloco de suavização iterativa de terreno da Fase~2
(\cref{sec:suavizacao}): o log real de produção revela que o binário do
\texttt{init\_atmosphere\_model} foi compilado com \texttt{"Default real
precision: single"} -- isto é, precisão simples (\texttt{real*4}) em toda
a base de código, não apenas nesse bloco. Como o critério de aceitação
por limiar desse algoritmo (\cref{eq:suavizacao_hx}) é sensível a
diferenças mínimas de arredondamento, compostas ao longo de até 85
passadas por até 53 níveis, as variáveis diretamente envolvidas nessa
suavização (terreno, $h_x$, acumuladores) foram declaradas com um
parâmetro de precisão simples dedicado
(\texttt{SPKIND=selected\_real\_kind(6)}), propositalmente, para casar
com o binário real -- uma escolha investigada e confirmada
experimentalmente (\cref{achado:hybrid}) como \emph{não} sendo, sozinha,
a causa da divergência numérica original, mas mantida por rigor mesmo
após a causa real ter sido encontrada, já que reduz ruído residual de
comparação.

\section{Leitura do \textit{namelist} real em tempo de execução}
\label{sec:namelist_config}

O módulo \texttt{namelist\_config.F90} define um tipo derivado
(\texttt{init\_atm\_config\_type}) que espelha os campos de configuração
efetivamente usados pelas fases deste trabalho, com valores-padrão
sensatos, e uma rotina (\texttt{read\_init\_atm\_namelist}) que lê, via
bloco \texttt{NAMELIST} nativo do Fortran, os grupos
\texttt{\&nhyd\_model}, \texttt{\&dimensions}, \texttt{\&data\_sources},
\texttt{\&vertical\_grid}, \texttt{\&interpolation\_control} e
\texttt{\&preproc\_stages} do arquivo \texttt{namelist.init\_atmosphere}
real do experimento -- o mesmo arquivo já consumido pela rota original,
sem nenhum formato novo introduzido. A leitura é tolerante a grupos ou
variáveis ausentes de um \textit{namelist} específico (por exemplo,
\texttt{config\_fg\_interval} só existe no \textit{namelist} do
\texttt{lbc\_run}, não no do \texttt{init\_run}): os valores-padrão do
tipo derivado são inicializados antes da leitura, de modo que uma
variável ausente simplesmente mantém seu padrão em vez de ficar
indefinida.

\begin{lstlisting}[caption={Trecho ilustrativo de \texttt{namelist\_config.F90}: leitura tolerante de múltiplos grupos de \textit{namelist}.}, label={lst:namelist}]
type :: init_atm_config_type
    real (kind=RKIND) :: config_ztop = 30000.0_RKIND
    integer            :: config_nsm  = 30
    real (kind=RKIND) :: config_dzmin = 0.3_RKIND
    logical            :: config_blend_bdy_terrain = .false.
    logical            :: config_frac_seaice       = .true.
    character (len=64) :: config_start_time = ''
    ! ... demais campos config_*
end type init_atm_config_type

subroutine read_init_atm_namelist(filename, cfg)
    ! ...
    open(newunit=unit_nml, file=trim(filename), status='old', action='read')
    read(unit_nml, nml=nhyd_model,    iostat=ios); rewind(unit_nml)
    read(unit_nml, nml=vertical_grid, iostat=ios); rewind(unit_nml)
    ! ... demais grupos, cada um tolerante a ausencia
    close(unit_nml)
    call parse_config_start_time(cfg)
end subroutine
\end{lstlisting}

Essa generalização, aplicada de forma consistente aos dois programas
principais (\texttt{gen\_init\_native}, \texttt{gen\_lbc\_native}),
resultou em uma constatação importante de validação: recompilar e
re-executar toda a cadeia com os parâmetros agora lidos dinamicamente,
em vez de fixos, produziu um resultado \textbf{numericamente idêntico,
bit a bit}, ao obtido com os valores anteriormente fixos no
código-fonte -- confirmando que a generalização não alterou nenhum
cálculo, apenas a origem dos parâmetros.

\section{Reúso de dependências e organização do build}

O sistema de \textit{build} (\texttt{Makefile}) mantém a convenção já
estabelecida no projeto \selfgrib{} de compilar \texttt{convert\_mpas}
como dependência externa antes do restante da pipeline, já que
\texttt{hinterp\_native} (Fase~1) reutiliza diretamente seus módulos
objeto já compilados (\texttt{mpas\_mesh.o}, \texttt{target\_mesh.o},
\texttt{remapper.o}, \texttt{scan\_input.o}), sem recompilá-los. Os
demais módulos novos (Fases~2 a~7) dependem apenas de
\texttt{mpas\_kind\_types.F90} e \texttt{mpas\_constants.F90}, ambos já
existentes no projeto -- nenhuma dependência externa nova foi introduzida
por este trabalho, além do utilitário NCO (\texttt{ncks}) usado na
composição final do \texttt{init.nc} (\cref{cap:fase7}), tipicamente já
disponível em ambientes de HPC voltados a geociências.
